\clearpage

\section{실험 및 성능 평가}

본 절에서는 제안된 AST 기반 GNN 모델의 취약점 탐지 성능을 실제 데이터셋을 통해 평가한 결과를 제시한다.

\subsection{실험 환경}

\subsubsection{데이터셋}
실험에는 NIST Juliet Test Suite v1.3의 C/C++ 데이터셋을 활용하였다. Juliet Test Suite는 다양한 소프트웨어 취약점 유형을 체계적으로 포함하는 표준 벤치마크 데이터셋으로, 본 실험에서는 다음 취약점 유형을 포함한다.

\begin{itemize}
    \item \textbf{CWE-121}: Stack-based Buffer Overflow
    \item \textbf{CWE-122}: Heap-based Buffer Overflow
    \item \textbf{CWE-124}: Buffer Underwrite
    \item \textbf{CWE-126}: Buffer Over-read
\end{itemize}

각 취약점 유형은 \texttt{bad} (취약) 및 \texttt{good} (안전/패치됨) 두 가지 변형을 포함하며, 다양한 코드 패턴과 문맥을 제공한다.

\subsubsection{실험 설정}
모델 학습 및 평가는 다음과 같은 설정으로 수행되었다.

\begin{itemize}
    \item \textbf{데이터 분할}: 전체 데이터셋을 학습 세트와 테스트 세트로 분할 (비율은 데이터셋 특성에 따라 결정)
    \item \textbf{모델 아키텍처}: GINEStack 기반 인코더, 히든 차원 128, 레이어 수 3
    \item \textbf{학습 파라미터}: 학습률 0.001, 배치 크기 32, 에포크 수 50
    \item \textbf{하드웨어}: NVIDIA GPU를 활용한 병렬 학습
\end{itemize}

\subsection{성능 지표}

\subsubsection{전체 성능}
테스트 세트 4,218개 샘플에 대한 평가 결과는 다음과 같다.

\begin{table}[h]
    \centering
    \begin{tabular}{lcc}
        \toprule
        \textbf{지표} & \textbf{값} & \textbf{비율} \\
        \midrule
        총 샘플 수 & 4,218 & 100.00\% \\
        정확도 (Accuracy) & 3,996 & 94.74\% \\
        \bottomrule
    \end{tabular}
    \caption{전체 성능 지표}
    \label{tab:overall_performance}
\end{table}

\subsubsection{혼동 행렬}
이진 분류 결과를 혼동 행렬(Confusion Matrix)로 정리하면 Table~\ref{tab:confusion_matrix}와 같다.

\begin{table}[h]
    \centering
    \begin{tabular}{l|cc|c}
        \toprule
        \multirow{2}{*}{\textbf{실제}} & \multicolumn{2}{c|}{\textbf{예측}} & \multirow{2}{*}{\textbf{합계}} \\
        \cmidrule{2-3}
        & 취약 (1) & 안전 (0) & \\
        \midrule
        취약 (1) & 3,441 (TP) & 129 (FN) & 3,570 \\
        안전 (0) & 93 (FP) & 555 (TN) & 648 \\
        \midrule
        \textbf{합계} & 3,534 & 684 & 4,218 \\
        \bottomrule
    \end{tabular}
    \caption{혼동 행렬 (Confusion Matrix)}
    \label{tab:confusion_matrix}
\end{table}

\subsubsection{상세 성능 지표}
혼동 행렬을 바탕으로 계산된 주요 성능 지표는 Table~\ref{tab:detailed_metrics}와 같다.

\begin{table}[h]
    \centering
    \begin{tabular}{lc}
        \toprule
        \textbf{지표} & \textbf{값} \\
        \midrule
        정확도 (Accuracy) & 94.74\% \\
        재현율 (Recall / TPR) & 96.39\% \\
        정밀도 (Precision) & 97.37\% \\
        F1-Score & 96.88\% \\
        오탐율 (FPR) & 14.35\% \\
        특이도 (Specificity) & 85.65\% \\
        \bottomrule
    \end{tabular}
    \caption{상세 성능 지표}
    \label{tab:detailed_metrics}
\end{table}

\subsection{성능 분석}

\subsubsection{높은 재현율의 의미}
재현율(Recall) 96.39\%는 실제 취약점 중 96.39\%를 정확히 탐지함을 의미한다. 이는 보안 도구로서 가장 중요한 지표 중 하나로, False Negative(미탐)를 최소화하여 실제 취약점을 놓치지 않는 것이 핵심이기 때문이다. 본 모델은 3,570개의 실제 취약점 중 3,441개를 정확히 탐지하여, 단 129개(3.61\%)만을 놓쳤다.

\subsubsection{오탐율 분석}
오탐율(False Positive Rate) 14.35\%는 정상 코드 648개 중 93개를 잘못 취약점으로 분류한 것이다. 이는 기존 정적 분석 도구의 전형적인 오탐율(20-40\%) 대비 경쟁력 있는 수치이나, 여전히 개선의 여지가 있다. False Positive를 감소시키기 위해서는 다음과 같은 방향의 개선이 필요하다.

\begin{itemize}
    \item \textbf{데이터 보강}: 정상 코드 패턴의 다양성을 높인 학습 데이터 추가
    \item \textbf{특징 개선}: 문맥 정보를 더 정교하게 반영하는 노드 특징 벡터 설계
    \item \textbf{앙상블 기법}: 여러 모델의 예측을 결합하여 오탐 감소
\end{itemize}

\subsubsection{클래스 불균형 대응}
데이터셋은 심각한 클래스 불균형을 보인다 (취약: 3,570개, 안전: 648개, 비율 약 5.5:1). 이러한 불균형 상황에서도 94.74\%의 높은 정확도를 달성한 것은 다음과 같은 요인에 기인한다.

\begin{itemize}
    \item \textbf{가중치 조정}: 학습 시 클래스 가중치(Class Weight)를 적용하여 소수 클래스(안전)에 더 높은 학습 가중치 부여
    \item \textbf{특징 품질}: AST 구조 정보가 취약점과 정상 코드를 효과적으로 구분할 수 있는 특징을 제공
    \item \textbf{모델 용량}: GINEStack의 충분한 표현력이 복잡한 패턴 학습을 가능하게 함
\end{itemize}

\subsection{오분류 분석}

\subsubsection{높은 확신도의 오분류}
모델 평가 결과, 일부 샘플에 대해서는 높은 확신도(Confidence)를 가지고 오분류하는 사례가 발견되었다. 예를 들어, \texttt{CWE122\_...\_rand\_12\_bad.json} 파일의 경우 모델이 0.99의 높은 확신도로 취약점으로 분류했으나 실제로는 정상 코드였다.

이러한 현상은 다음과 같은 원인으로 분석된다.

\begin{itemize}
    \item \textbf{패턴 유사성}: 일부 정상 코드가 취약점 코드와 매우 유사한 구조적 패턴을 가짐
    \item \textbf{문맥 정보 부족}: AST 구조만으로는 판단하기 어려운 런타임 검증 로직의 부재
    \item \textbf{데이터 편향}: 학습 데이터에 특정 패턴이 과도하게 포함되어 모델이 이를 과도하게 일반화
\end{itemize}

\subsubsection{개선 방향}
높은 확신도의 오분류를 줄이기 위한 향후 연구 방향은 다음과 같다.

\begin{itemize}
    \item \textbf{Adversarial Training}: 오분류된 샘플을 학습에 포함시켜 모델의 견고성 향상
    \item \textbf{데이터 보강}: 오분류 패턴을 포함하는 추가 학습 데이터 수집 및 생성
    \item \textbf{확신도 기반 필터링}: 낮은 확신도 예측에 대해서는 추가 검증 절차 적용
    \item \textbf{설명 가능성 분석}: 오분류 원인을 분석하여 모델의 의사결정 과정 개선
\end{itemize}

\subsection{결론}

본 실험을 통해 제안된 AST 기반 GNN 모델이 실제 취약점 탐지 작업에서 매우 높은 성능을 보임을 입증하였다. 특히 96.39\%의 높은 재현율은 보안 도구로서 실용적 가치가 높음을 시사한다. 다만, 14.35\%의 오탐율은 추가적인 개선이 필요하며, 향후 데이터 보강 및 모델 개선을 통해 더욱 정확한 탐지가 가능할 것으로 기대된다.



